\documentclass[12pt]{ctexart}
\usepackage{tikz}
\usepackage{pgfplots}
\usepackage{amsmath}
\pgfplotsset{compat=1.18}

\begin{document}
% 1. 等效直径
% grid=none        % 无网格（默认）
% grid=major       % 只显示主刻度网格（最常用）
% grid=minor       % 只显示副刻度网格
% grid=both        % 主+副网格都显示
\begin{tikzpicture}
    \begin{axis}[
        title={等效直径 $d = \sqrt{\dfrac{4A}{\pi}}$},
        xlabel={面积 $A$},
        ylabel={直径 $d$},
        grid=major,
    ]
    \addplot[blue, thick, domain=0:100]{sqrt(4*x/pi)};
    \end{axis}
\end{tikzpicture}

% 2. 函数曲线图
\begin{tikzpicture}
    \begin{axis}[
        title={函数曲线图 $y=x^3$},  
        xlabel={$x$},
        ylabel={$y$},
        grid=major,
    ]
    \addplot[blue, thick, domain=-2:2]{x^3};
    \end{axis}
\end{tikzpicture}

% 3. 散点图
\begin{tikzpicture}
    \begin{axis}[
        title={散点图},      
        xlabel={X},
        ylabel={Y},
        grid=major,
    ]
    \addplot[only marks, red, mark=o] coordinates {
        (1, 2)
        (2, 3)
        (3, 1)
        (4, 5)
    };
    \end{axis}
\end{tikzpicture}

% 4. 柱状图
\begin{tikzpicture}
    \begin{axis}[
        title={柱状图},      
        ybar,
        xtick={1,2,3,4},
        % xlabel={X}
    ]
    \addplot coordinates {
        (1, 5)
        (2, 3)
        (3, 8)
        (4, 6)
    };
    \end{axis}
\end{tikzpicture}

% 5. 多条曲线
% red        红
% green      绿
% blue       蓝
% yellow     黄
% cyan       青
% magenta    品红
% black      黑
% white      白
% gray       灰
% darkgray   深灰
% lightgray  浅灰
% orange     橙
% violet     紫
% purple     紫
% brown      棕
% pink       粉
\begin{tikzpicture}
    \begin{axis}[
        title={多条曲线},    
        grid=both,
        domain=-2*pi:2*pi,
        samples=100,    % 采样点数默认是 25
        legend pos=north west,
    ]
    \addplot[red]{x};
    \addlegendentry{$y=x$}
    \addplot[blue]{x^2};
    \addlegendentry{$y=x^2$}
    \addplot[brown]{10*sin(deg(x))};
    \addlegendentry{$y=\sin x$}
    \end{axis}
\end{tikzpicture}

\end{document}